\documentclass[10pt]{article}
\usepackage[margin=0.8in]{geometry}
\usepackage{fontspec}
\setmainfont{DejaVu Serif}
\setmonofont{DejaVu Sans Mono}
\usepackage{booktabs}
\usepackage{longtable}
\usepackage{array}
\usepackage{hyperref}
\usepackage{xcolor}
\usepackage{listings}
\usepackage{enumitem}
\usepackage{fancyhdr}
\hypersetup{colorlinks=true,linkcolor=blue,urlcolor=blue}
\setlength{\parskip}{0.45em}
\setlength{\parindent}{0pt}
\pagestyle{fancy}
\fancyhf{}
\lhead{RunPod historical incident report}
\rhead{2026-08-06}
\cfoot{\thepage}
\lstset{basicstyle=\ttfamily\scriptsize,breaklines=true,columns=fullflexible,frame=single}
\newcommand{\obs}{\textbf{Observed.} }
\newcommand{\inferr}{\textbf{Inferred.} }
\newcommand{\decision}{\textbf{Decision.} }
\begin{document}

\begin{center}
{\LARGE RunPod attempts, outcomes, and workflow incident review}\\[0.4em]
{\large Prepared for Benjamin Goertzel --- 2026-08-06}\\[0.6em]
\textit{Scope: documented attempts visible in the research workspace and current-session operational records.}
\end{center}

\section*{Executive finding}

RunPod is demonstrably usable from this account and machine.  The historical
record includes successful authenticated SSH canaries, successful raw-image
and template launches, completed GPU experiments with retrieved artifacts,
and prompt cleanup.  The recent RelaLeap C1 sequence nevertheless suffered
from a compound agent-process failure: lifecycle ownership was not kept single
writer; stale monitors produced contradictory reports; provider control-plane
``RUNNING'' was mistaken for usable runtime; and a pre-update CLI was used
without its now-available bounded readiness wait.

This does not establish that every recent non-ready allocation was caused by
the agent.  Some allocations were never ready or disappeared provider-side.
But the agent's workflow made those events costly, poorly observable, and
confusing.  It was not an adequate procedure for a popular, normally usable
service.

\section*{Method, provenance, and limits}

This review read the durable project records, experiment ledgers, returned
artifact records, relevant daily-memory entries, the current RunPod CLI help,
and the exact current \texttt{runpodctl} skill reproduced in Appendix A. It
does not claim access to RunPod's internal scheduler, image-pull logs, or
provider billing ledger beyond captured CLI outputs.  A pod's requested status
is therefore reported as a control-plane observation, not as proof that a
container ran.

Sensitive material was excluded: API keys, SSH private keys, relay codes, and
full environment dumps are not reproduced.  Pod IDs are included because they
are essential provenance, but no credential-bearing command output appears.

\section*{Outcome taxonomy}

\begin{itemize}[leftmargin=1.2em]
\item \textbf{Operational success}: remote runtime was reached and the intended
  narrow operation completed or decisively demonstrated a capability.
\item \textbf{Scientific negative result}: remote execution and artifact
  recovery succeeded, but the research hypothesis failed its frozen gate.
\item \textbf{Code/launcher failure}: provider runtime worked, but project
  software failed; this must not be attributed to RunPod.
\item \textbf{Provider-readiness failure}: an allocation did not become a
  usable runtime, or disappeared, before the intended work.
\item \textbf{Agent workflow failure}: authority, monitoring, observability,
  or lifecycle procedure was itself defective.
\end{itemize}

\section*{Evidence-backed timeline}

\small
\begin{longtable}{p{0.11\textwidth}p{0.22\textwidth}p{0.23\textwidth}p{0.31\textwidth}}
\toprule
Date & Attempt & Outcome & What the evidence establishes \\
\midrule
\endhead
2026-07-10/11 & RelaLeap HDPC grid, pod \texttt{zivpxu06yhmx40}, \$0.49/hr & Operational success & GPU grid completed; 9 result entries and 9 checkpoints were retrieved, JSON was parsed, and the pod was deleted. This proves a normal GPU workload, return path, and cleanup worked. [S1] \\
2026-07-11 & RelaLeap SLT large run, pod \texttt{pmqngq3cxcd5bw}, A100 \$1.49/hr & Operational success & CUDA result, 8 chains x 1000 samples, and a 3978 s run were retrieved and checked before termination. Scientific convergence caveats do not negate operational success. [S2] \\
2026-07-17 & RelaLeap A100 outcome attempt, pod \texttt{lop52u56xugpxk} & Provider-readiness failure & Allocation remained unready with zero uptime/SSH for about 52 minutes and was terminated. This is a genuine historical non-ready class, but not an explanation for every later event. [S3] \\
2026-07-18/19 & RelaLeap six-layer runs & Mixed: code failure, then successful negative experiment & Run 4 exposed application bugs (tensor shape and local checkpoint path). The r5 rerun, pod \texttt{oo20lk2075y0dx}, completed all three seeds and returned 13 JSON artifacts; its promotion gate failed scientifically. Provider use itself succeeded. [S4] \\
2026-07-29 & SSH-free Croc transport probe & Agent-design/observability failure & Three A40 pods returned no probe artifact; the receiver's ``room not ready'' could not establish whether the remote startup ran. The approved envelope was exceeded. Later review identified relay-code/docker-args misuse; this was not evidence that RunPod itself failed. [S5] \\
2026-07-29 & Standard RTX 4090 sweep, pod \texttt{6z14ahtjar5bzh} & Runtime success; launcher contract failure & SSH and GPU execution were observed. The run failed because the runner emitted an objective tuple while the validator required a list; artifacts were retrieved and hash-verified. [S6] \\
2026-07-30 & Clean-room robustness run & Runtime success; project gate failure & CUDA/bootstrap passed; then an explicit \texttt{trust\_remote\_code=True} dependency caused a fail-closed project stop before seed execution. Evidence was retrieved. [S7] \\
2026-07-30/08-01 & Clean-room confirmation and reliability runs & Operational success; mixed scientific results & One three-seed confirmation completed with positive frozen result; the later reliability run completed and failed its primary scientific gate, with returned hashes verified and pod deleted. [S7] \\
2026-07-31 & Template SSH canary, pod \texttt{50tq84z1qxiok3} & Operational control success & Official template, explicit \texttt{22/tcp}, and managed-key SSH reached a Secure RTX 4090 in 47 s, then deleted in 74 s. This rules out an account-wide SSH/key/API failure at that time. [S8] \\
2026-07-31 & Raw-image A/B, pod \texttt{tdhmwpcs5v77r9} & Operational control success & Same image via \texttt{--image}, explicit ports/SSH, reached authenticated SSH in 34 s; PID 1 included \texttt{/start.sh}. Raw images are not intrinsically incompatible with this account and CLI. [S9] \\
2026-08-05/06 & C1 A40 sequence, IDs \texttt{dnwq8koj3rtu45}, \texttt{elatgmvrthhxf7}, \texttt{20byqukvf2kzv4} & Mixed provider and agent workflow failure & First two were created outside a stable single-writer envelope and died before readiness. The third did reach SSH/GPU, but was bound to a mutable tag while the sealed handoff required a digest, so no science ran; it later disappeared. Stale monitors continued reporting the dead second pod. [S10] \\
2026-08-06 & C1 RTX 6000 Ada, pod \texttt{crgb1jzlu28uix} & Provider-readiness failure amplified by agent procedure & It remained requested \texttt{RUNNING} with uptime 0 and ``pod not ready'' for about 95 min; it was then deleted. The agent used CLI 2.8.1 without its required update and did not gate creation with a bounded readiness wait. [S11] \\
\bottomrule
\end{longtable}
\normalsize

\section*{What succeeded versus what failed}

The workspace contains multiple successful direct-SSH and GPU-run examples.
Therefore the statement ``the account cannot use RunPod'' is false.  It also
contains genuine non-ready allocations and disappearing provider objects. The
more accurate statement is: the provider can sometimes fail readiness, while
the recent agent procedure failed to detect, contain, and explain that state
promptly.

\begin{center}
\begin{tabular}{p{0.29\textwidth}p{0.58\textwidth}}
\toprule
Claim & Status \\
\midrule
Account/API/managed SSH universally broken & Refuted by July 31 template and raw-image controls. \\
Raw \texttt{--image} intrinsically prevents startup/SSH & Refuted by the raw-image A/B control. \\
Every non-ready C1 pod was an agent bug & Not established; provider readiness failures are observed. \\
Recent C1 lifecycle procedure was adequate & Refuted by concurrent writers, stale monitors, unbounded non-ready billing, and mutable-image mismatch. \\
Secure Cloud means on-demand reliability & Unsupported; the current CLI exposes \texttt{SECURE} versus \texttt{COMMUNITY} as cloud type, not a reliability guarantee. \\
\bottomrule
\end{tabular}
\end{center}

\section*{Root-cause analysis}

\subsection*{Provider-side uncertainty}

\obs Several allocations reported zero uptime or no SSH; some later vanished
from both active and historical lists.  The local CLI did not supply internal
host, image-pull, or scheduler diagnostics.  The evidence supports a
provider-readiness class, but cannot distinguish placement failure, host
initialization, image pull, or provider-side deletion.

\inferr An image digest may be relevant to an individual failure, but it is
not proven here.  The proper next experiment is a bounded current-CLI canary
using a known-good official template (and optionally a controlled raw-image
comparison), not more C1 launches.

\subsection*{Agent-side failures}

\begin{enumerate}[leftmargin=1.5em]
\item \textbf{Readiness was not a hard admission gate.} Desired \texttt{RUNNING}
  was treated as progress even when uptime was zero and no SSH banner existed.
\item \textbf{The CLI was not updated before paid use.} The installed skill
  explicitly said to update/pin a recent CLI; the C1 RTX 6000 Ada attempt used
  2.8.1.  After update, 2.9.0 exposes \texttt{pod create --wait --wait-timeout},
  which directly encodes the missing readiness gate.
\item \textbf{Lifecycle ownership was not exclusive.} Concurrent writers
  created out-of-envelope C1 pods; later monitors retained obsolete pod IDs.
\item \textbf{Routine monitoring was noisy and semantically stale.} It
  reported ``no material change'' repeatedly despite instructions not to do
  so, wasting attention and masking the relevant state.
\item \textbf{Cloud type was overinterpreted.} \texttt{SECURE} was called
  ``on-demand'' without evidence that it conveyed the requested persistence
  property.
\item \textbf{C1 sealing was not prepared at the correct lifecycle boundary.}
  One reachable pod was correctly prevented from running because it advertised
  a mutable tag while the sealed handoff expected a digest. That fail-closed
  behavior was correct; provisioning it before resolving the bundle identity
  was not.
\end{enumerate}

\section*{Corrected minimal RunPod protocol}

No new paid attempt is authorized by this report.  Before the next one, the
following protocol should be followed and recorded:

\begin{enumerate}[leftmargin=1.5em]
\item Update and pin \texttt{runpodctl}; record \texttt{version} and a healthy
  \texttt{doctor} result.
\item Create one explicit remote-job record and elect one lifecycle owner.
  Disable or remove every monitor for superseded pod IDs.
\item Use an official known-good template with explicit \texttt{22/tcp} and
  SSH for a bounded canary.  Create with \texttt{--wait --wait-timeout 10m}.
\item Treat only a successful SSH banner plus authenticated command and GPU
  identity check as ready.  On timeout, delete the exact returned ID and query
  the provider again; do not wait indefinitely or create a replacement.
\item Only after a canary passes, create the scientific pod.  Pin image and
  source digests before provisioning; bind the sealed handoff to the actual
  resource identity; transfer minimum non-sensitive files.
\item Use one silent monitor that emits only state transitions, hard-stop,
  success, or failure.  Retrieve/hash artifacts before termination and record
  the final provider-empty check.
\end{enumerate}

\section*{Recommended controlled validation}

The next useful validation is not another C1 run.  It is a single current-CLI
canary under a fresh, small cost envelope: known-good official template,
Secure RTX 4090 or comparable broadly available GPU, 20 GB ephemeral disk,
explicit ports/SSH, \texttt{--wait --wait-timeout 10m}, one benign remote
identity command, and immediate termination.  Its acceptance criterion is
not ``RUNNING''; it is successful authenticated SSH plus \texttt{nvidia-smi}
and a post-delete empty provider query.  A failed canary should preserve the
exact create response and timeout ID, then stop rather than retry automatically.

\section*{Source register}

\begin{description}[style=nextline,leftmargin=1.2em]
\item[S1] \texttt{projects/relaleap/experiments/hdpc\_epc\_grid\_run1/RUN.md}.
\item[S2] \texttt{memory/2026-07-11.md}, ``RelaLeap SLT large Runpod cleanup''.
\item[S3] \texttt{projects/relaleap/TASKS.md}, July 17 A100 readiness entry.
\item[S4] \texttt{projects/relaleap/experiments/r5-runpod-20260719/EXPERIMENT.md};
\texttt{projects/relaleap/TASKS.md}, July 18--19 records.
\item[S5] \texttt{projects/remote-job-bootstrap/experiments/20260729T1555Z-remote-transport-probe/RUN.md}.
\item[S6] \texttt{memory/2026-07-29.md}, ``RunPod bootstrap / RelaLeap''.
\item[S7] \texttt{projects/relaleap/PROJECT.md}, clean-room July 29--August 1 entries.
\item[S8] \texttt{projects/remote-job-bootstrap/experiments/20260731T133317Z-template-ssh-canary/RUN.md}.
\item[S9] \texttt{projects/remote-job-bootstrap/experiments/20260731T140535Z-raw-image-ssh-ab/RUN.md}.
\item[S10] \texttt{projects/relaleap/experiments/20260806T0026Z-epc-c1-any-gpu-runpod/REMOTE\_JOB.md};
\texttt{projects/relaleap/DECISIONS.md}; \texttt{memory/2026-08-06.md}.
\item[S11] \texttt{projects/relaleap/experiments/20260806T133500Z-epc-c1-secure-rtx6000ada/REMOTE\_JOB.md};
current-session CLI receipts from 2026-08-06.
\end{description}

\clearpage
\appendix
\section*{Appendix A: Current RunPod skill (verbatim)}

Source path at report generation:\\
\texttt{/home/openclaw/.openclaw/agents/main/agent/codex-home/skills/runpodctl/SKILL.md}\\
SHA-256: \texttt{ae4850e1e297be9bc46f443f0e7b46d68ce84844ba262745638602b80047081a}\\
Line count: 248.

\lstinputlisting{/home/openclaw/.openclaw/agents/main/agent/codex-home/skills/runpodctl/SKILL.md}

\end{document}
